% =====================================================================
\section{Notation, Conventions and Parameters}
\label{sec:notation}
% =====================================================================

\subsection{Frames and conventions}

\begin{itemize}
  \item \textbf{Inertial frame} $\mathcal F_i$ (in Gazebo: the world frame, $z$ up) and \textbf{body frame} $\mathcal F_b$,
        fixed to the vehicle at its centre of mass. The angular velocity $\vw$ is the rate of $\mathcal F_b$ with
        respect to $\mathcal F_i$, \emph{expressed in} $\mathcal F_b$ (what a gyroscope measures).
  \item \textbf{Quaternion order.} The paper writes $\vq=[q_1,q_2,q_3,q_4]^\top$ with the scalar \emph{last}. ROS and
        Gazebo use $[x,y,z,w]$, also scalar last, so the two agree and no reordering is needed. (In its Lyapunov
        function the paper calls the scalar part $q_0$; it is the same $q_4$.)
  \item \textbf{Element-wise operations.} Inside the barrier term the division and the power act on each axis
        separately: $(1-\veps_1)^{-2}$ means the vector with entries $(1-\varepsilon_{1i})^{-2}$, and the bracket multiplies
        $\vs$ entry by entry (a diagonal gain matrix).
\end{itemize}

\subsection{Master symbol table}

\begin{center}\small
\begin{longtable}{@{}>{$}l<{$} p{5.6cm} >{$}l<{$} l p{3.3cm}@{}}
\toprule
\textrm{\textbf{Symbol}} & \textbf{Meaning} & \textrm{\textbf{Dim.}} & \textbf{Unit} & \textbf{Paper value}\\
\midrule\endhead
\vq=[\qv^\top,q_4]^\top & unit attitude quaternion (error to the target) & \R^4 & -- & $\|\vq\|=1$\\
\qv & vector part of $\vq$ & \R^3 & -- & --\\
\vw & body angular velocity & \R^3 & rad/s & --\\
J & inertia matrix, symmetric, $J\succ0$ & \R^{3\times3} & kg\,m$^2$ & see \eqref{eq:Jpaper}\\
\vu & torque commanded by the controller & \R^3 & N\,m & --\\
\sat(\vu) & torque the actuators really produce & \R^3 & N\,m & $|\cdot|\le\umax$\\
\umax & actuator limit per axis & \R & N\,m & 2.5\\
\vd & external disturbance torque, $\|\vd\|\le\bar d$ & \R^3 & N\,m & see \eqref{eq:dpaper}\\
\bm\Pi=[\pi_1,\pi_2,\pi_3]^\top & torque lost to saturation, $\sat(\vu)-\vu$ & \R^3 & N\,m & --\\
\bm T_d=\vd+\bm\Pi & lumped unknown torque & \R^3 & N\,m & --\\
\theta & unknown bound, $\|\bm T_d\|\le\theta$ & \R & N\,m & --\\
\that,\ \ttil=\theta-\that & estimate of $\theta$ and its error & \R & N\,m & $\that(0)=0$\\
\rho_q(t),\ \rho_\omega(t) & envelopes (tube half-widths) for $\qv$ and $\vw$ & \R & --, rad/s & --\\
\rho_{q0},\rho_{qT} & initial and final width of $\rho_q$ & \R & -- & 0.3, 0.008\\
\rho_{\omega0},\rho_{\omega T} & initial and final width of $\rho_\omega$ & \R & rad/s & 0.4, 0.005\\
\kappa & shape parameter of the FTPPF (paper: $\sigma_q,\sigma_\omega$) & \R & -- & 2.5\\
T_s & settling time of the FTPPF (paper: $T_f$) & \R & s & 15 and 10\\
\veps_1,\ \veps_2 & normalised errors $\qv/\rho_q$, $\vw/\rho_\omega$ & \R^3 & -- & $|\cdot|<1$\\
\vs=\vw+c\,\qv & sliding (composite) variable & \R^3 & rad/s & --\\
c & weight of $\qv$ in $\vs$ & \R & 1/s & \emph{not stated}; we use 0.5\\
k_p,\ k_d & PD gains & \R & N\,m, N\,m\,s & 40, 20\\
\alpha,\ \beta & barrier weights & \R & N\,m\,s & 1.5, 1.2\\
\vu_{cons} & auxiliary (constraint) torque & \R^3 & N\,m & --\\
\varsigma & boundary-layer width of $\sat(s/\varsigma)$ & \R & rad/s & \emph{not stated}; we use 0.02\\
\eta_1 & adaptation time constant & \R & -- & 0.1\\
\eta_2(t) & leakage, $\int_0^\infty\eta_2\,dt<\infty$ & \R & 1/s & $2e^{-0.5t}$\\
\bottomrule
\end{longtable}
\end{center}

\subsection{The paper's numerical set-up (Section IV)}

\begin{equation}\label{eq:Jpaper}
J=\begin{bmatrix}20&1.2&0.9\\1.2&17&1.4\\0.9&1.4&15\end{bmatrix}\ \text{kg\,m}^2,\qquad
\lambda(J)=\{14.27,\ 17.00,\ 20.74\}\ \text{kg\,m}^2
\end{equation}
\begin{equation}\label{eq:dpaper}
\vd(t)=\big[\,0.05\sin t,\ \ 0.1\sin 1.2t,\ \ 0.15\sin 1.5t\,\big]^\top\ \text{N\,m},\qquad \|\vd\|\le0.187\ \text{N\,m}
\end{equation}

\begin{center}\small
\begin{tabularx}{\linewidth}{@{}l X X@{}}
\toprule
& \textbf{Part 1} (effect of $T_s$) & \textbf{Part 2} (constrained vs.\ conventional PD)\\
\midrule
$\vq(0)$ & $[0.2,\ -0.15,\ -0.25,\ 0.9354]^\top$ & $[0.3,\ -0.2,\ -0.3,\ 0.8832]^\top$\\
$\vw(0)$ (rad/s) & $[0.03,\ 0.02,\ -0.01]^\top$ & $[-0.02,\ -0.02,\ 0.025]^\top$\\
$T_s$ & 15\,s and 10\,s & not stated; we use 15\,s\\
$\rho_{q0}$ & 0.3 & 0.3 in the paper; we use 0.36 (see \cref{warn:part2})\\
\bottomrule
\end{tabularx}
\end{center}

\begin{warn}[Four values the paper does not give, and what we chose]
\label{warn:part2}
\begin{enumerate}[leftmargin=*]
  \item $c$ in $\vs=\vw+c\qv$. Theorem~1 needs $c$ ``sufficiently small''; the explicit condition (\cref{sec:theorem}) is
        $c<k_d/\lambda_{\max}(J)=20/20.74=0.964$. We use $c=0.5$.
  \item $\varsigma$ in $\sat(s/\varsigma)$. We use $\varsigma=0.02$\,rad/s, small compared with the initial $\|\vs\|\approx0.2$.
  \item The settling time of Part~2. We use $T_s=15$\,s, the first value of Part~1.
  \item $\rho_{q0}$ in Part~2. The paper's Part~2 initial attitude has $|q_1(0)|=|q_3(0)|=0.3=\rho_{q0}$, so
        $\varepsilon_{11}(0)=1$: the state starts \emph{on} the boundary, the barrier $(1-\varepsilon)^{-2}$ is infinite and
        the paper's own requirement $|q_i(0)|<\rho_{q0}$ (Remark~7) is violated. We use $\rho_{q0}=1.2\max_i|q_i(0)|=0.36$.
\end{enumerate}
Every result in this guide and in the complete documentation is reported with these choices, and the sensitivity to
them is measured (complete documentation, Section ``Sensitivity'').
\end{warn}

\subsection{The X3 quadcopter (Gazebo) parameters}

All values are read from the Gazebo model files (\code{quadcopter.sdf} and the \code{MulticopterMotorModel} plugin), not
tuned. Rotor positions are in the body frame, $x$ forward, $y$ left.

\begin{center}\small
\begin{tabularx}{\linewidth}{@{}l X l@{}}
\toprule
\textbf{Quantity} & \textbf{Value} & \textbf{Source}\\
\midrule
mass $m$ & 1.5\,kg body + $4\times0.005$\,kg rotors $=1.52$\,kg & SDF\\
inertia $J$ about the CoM & $\diag(0.03565,\ 0.07035,\ 0.09892)$\,kg\,m$^2$ (body + rotors, parallel-axis) & SDF\\
rotor 0, 1, 2, 3 & $(0.13,-0.22)$ ccw, $(-0.13,0.20)$ ccw, $(0.13,0.22)$ cw, $(-0.13,-0.20)$ cw (m) & SDF\\
thrust constant $k_T$ & $8.54858\times10^{-6}$\,N/(rad/s)$^2$, $F_i=k_T\Omega_i^2$ & plugin\\
moment constant $c_Q$ & 0.016\,m, yaw torque $-\mathrm{dir}_i\,c_QF_i$ & plugin\\
max rotor speed & 800\,rad/s, so $F_{\max}=5.471$\,N per rotor & plugin\\
rotor time constants & 12.5\,ms (up), 25\,ms (down) & plugin\\
pose rate & $\approx$60\,Hz on \code{/world/quadcopter/dynamic\_pose/info} & SceneBroadcaster\\
\bottomrule
\end{tabularx}
\end{center}

\textbf{Scale.} The paper's spacecraft has an inertia about \textbf{250 times} larger than the X3 and a torque limit about 7
times larger. The X3 can produce at most 1.46\,N\,m in roll, 0.91\,N\,m in pitch and only \textbf{0.106\,N\,m in yaw} while
hovering (\cref{sec:b10}). The paper's gains and envelopes therefore cannot be copied; they must be re-sized
(\cref{sec:b9}).
